% GENERATED by scripts/build_paper.py from appendix.template.tex; edit the template, never this file.
% SOURCE of the appendix (single-column submission only; scripts/build_paper.py fills it into appendix.tex and
% main_single.tex inputs it after the references). Same rules as body.template.tex: every number is a double-brace
% placeholder read from a result file, and every analysis added in the revision is labelled (exploratory).
% The body refers to these sections as appendix A1 to A7; main_single.tex numbers them A1, A2, ...

\section{Training settings}\label{app:A1}
All rollout-trained models of the F18 and five-split pools were trained with the settings in Table~\ref{tab:app-train}.
The optimiser is Adam; gradient clipping: none. The loss has 7 non-zero terms,
with weights state 1, sequence 1, bound 1, kl 0.1, extension 1, contact 1, termination 1. These settings were checked to be identical in the training records of all
25 parent models of the five-split pool. The recurrent network has two GRU layers of 256 units
and reads 32 past steps.

\begin{table}[h]
\centering
\caption{What differs between the rollout-trained parent, the teacher-forced models and the $N{=}1$ models. A blank means
the setting equals the parent's. Records compared: 3 teacher-forced and 3 $N{=}1$ models.}
\label{tab:app-train}
\footnotesize
\fittable{\begin{tabular}{@{}lccc@{}}
\toprule
Setting & Rollout-trained parent & Teacher-forced & $N{=}1$ \\
\midrule
Learning rate / weight decay / batch & $1\times10^{-4}$ / $1\times10^{-5}$ / 256 & & \\
Iterations & 2{,}500 & & \\
Loss target & own predictions fed back & recorded states fed in & own predictions \\
Steps predicted per window ($N$) & 8 & 8 & 1 \\
\bottomrule
\end{tabular}}
\end{table}

Other F18 models differ from the parent in one setting each. The longer-trained models ran 10{,}000
iterations (the teacher-forced ones as well). The ensembles have 5 members. The capacity-matched control has
124 hidden units. The likelihood model replaces the squared error by a Gaussian likelihood, and two
further models train on contaminated or duplicated data. The compressed copies are made from the finished parents without
further training: 16-bit and 8-bit weights, removal of the smallest weights over the whole network (10, 20, 30 and
50\%), and removal of units so that the hidden width is 192 or 128.

\section{What was known before the registration}\label{app:A2}
Section 2 of \texttt{PREREG.md} states what the author had read before any model was run through this folder's code.
For the growth rate $\lambda_1$ on the F18 pool, an earlier study found Spearman +0.85 [+0.70, +0.93]
with the 368-step error over 48 rollout-trained models, and +0.90 [+0.80, +0.94] over
57 rollout-trained and teacher-forced ones. On the five splits, $\lambda_1$ ranked 55 models at
+0.66 [+0.46, +0.79], and the 40 rollout-trained ones among them at
+0.21 [$-$0.13, +0.50] (not settled). The mean kinematic residual ranked 55 models at
+0.93 [+0.84, +0.97], much of it through the difference between training methods. Earlier studies also
showed that teacher forcing beats rollout training at one step but loses at 368 steps, and that moderate pruning
without fine-tuning loses skill at long horizons.
Open were the one-step and 8-step errors computed from every position of the held-out data, the correlation of any
short-horizon error with any long-horizon error, and the growth rate or residual of any compressed copy. The raw
arrays of the earlier growth-rate study were used for one consistency check only, and no statistic is taken from them.

\section{Horizon-matched and split-half checks}\label{app:A3}
Table~\ref{tab:app-horizon} gives the Spearman correlation with the 100-step error $E_{100}$ when the growth rate
and the kinematic residual are computed over a window of 8, 24, 100 or 368 steps (exploratory). The one-step and 8-step
errors are shown for comparison. Over 8 steps the growth rate's correlation is +0.21 on F18 and
$-$0.02 on X5, against +0.77 and +0.56 over 100 steps.
With both checks computed over 8 steps, the residual's lead over the 8-step error was $-$0.006
[$-$0.076, +0.063] on F18 and $-$0.073 [$-$0.120, $-$0.040] on X5; over 100 steps it was
+0.109 [+0.035, +0.205] on F18 and +0.013 [$-$0.003, +0.033] on X5 (all exploratory).

\begin{table}[h]
\centering
\caption{Spearman correlation of each check with $E_{100}$ by the window the check is computed over, for F18 (4 held-out
trajectories; all simulators, trained, compressed) and X5 (20 trajectories) (exploratory).}
\label{tab:app-horizon}
\footnotesize
\setlength{\tabcolsep}{4pt}
\fittable{\begin{tabular}{@{}lcccc@{}}
\toprule
Check (window) & F18, all & F18, trained & F18, compressed & X5 \\
\midrule
One-step error $E_1$ & +0.32 & $-$0.50 & +0.89 & +0.78 \\
8-step error $E_8$ & +0.84 & +0.63 & +0.96 & +0.95 \\
\midrule
$\lambda_1$ (8 steps) & +0.21 & +0.72 & $-$0.60 & $-$0.02 \\
$\lambda_1$ (24 steps) & +0.33 & +0.80 & $-$0.30 & +0.19 \\
$\lambda_1$ (100 steps) & +0.77 & +0.90 & +0.58 & +0.56 \\
$\lambda_1$ (368 steps) & +0.80 & +0.85 & +0.82 & +0.67 \\
\midrule
$\bar r_J$ (8 steps) & +0.83 & +0.60 & +0.90 & +0.88 \\
$\bar r_J$ (24 steps) & +0.88 & +0.70 & +0.95 & +0.93 \\
$\bar r_J$ (100 steps) & +0.94 & +0.83 & +0.97 & +0.96 \\
$\bar r_J$ (368 steps) & +0.95 & +0.91 & +0.94 & +0.96 \\
\bottomrule
\end{tabular}}
\end{table}

Table~\ref{tab:app-split} computes each check on one half of the trajectories and the target on the other half
(exploratory), beside the correlation when both use the same half. The largest drop is 0.04.

\begin{table}[h]
\centering
\caption{Split-half check (exploratory): Spearman correlation with $E_{100}$ when the check and the target come from
different halves of the held-out trajectories, with its 95\% interval, and when both come from the same half.}
\label{tab:app-split}
\footnotesize
\setlength{\tabcolsep}{4pt}
\fittable{\begin{tabular}{@{}lcccc@{}}
\toprule
 & F18, other half & F18, same half & X5, other half & X5, same half \\
\midrule
$E_1$ & +0.31 [+0.07, +0.51] & +0.32 & +0.77 [+0.68, +0.83] & +0.78 \\
$E_8$ & +0.77 [+0.64, +0.86] & +0.79 & +0.91 [+0.87, +0.93] & +0.94 \\
$\lambda_1$ & +0.73 [+0.64, +0.79] & +0.77 & +0.62 [+0.54, +0.68] & +0.62 \\
$\bar r_J$ & +0.91 [+0.86, +0.94] & +0.92 & +0.91 [+0.86, +0.93] & +0.94 \\
\bottomrule
\end{tabular}}
\end{table}

\section{Wider intervals, the five-split pool without split \{1,8\}, and Holm's correction}\label{app:A4}
All of this section is exploratory. The registered intervals resample the held-out trajectories. The wider intervals
also resample the parent models, in two levels, with 20{,}000 draws. They are 1.5 times as wide as the
registered ones (median; 1.3 for the correlations of Table~\ref{tab:checks} and 1.6 for
the registered differences; range 1.0 to 1.8). Of the 13 registered verdicts,
10 hold and 3 do not.

\begin{table}[h]
\centering
\caption{Table~\ref{tab:checks}'s correlations with $E_{100}$, with the registered interval (where one was computed) and
the two-level 95\% interval (exploratory); for X5 also without the split \{1, 8\}, which repeats F18's held-out pair (exploratory).}
\label{tab:app-two}
\footnotesize
\setlength{\tabcolsep}{3pt}
\fittable{\begin{tabular}{@{}llcccc@{}}
\toprule
Pool & Check & Registered & Registered 95\% & Two-level 95\% & Without \{1, 8\} \\
\midrule
F18, all & $E_1$ & +0.32 & [+0.07, +0.54] & [$-$0.18, +0.59] & \\
 & $E_8$ & +0.84 & [+0.71, +0.93] & [+0.56, +0.91] & \\
 & $\lambda_1$ & +0.80 & [+0.71, +0.85] & [+0.63, +0.87] & \\
 & $\bar r_J$ & +0.95 & [+0.91, +0.97] & [+0.86, +0.97] & \\
F18, trained & $E_1$ & $-$0.50 & [$-$0.71, $-$0.20] & [$-$0.69, $-$0.06] & \\
 & $E_8$ & +0.63 & [+0.36, +0.83] & [+0.19, +0.80] & \\
 & $\lambda_1$ & +0.85 & [+0.72, +0.91] & [+0.70, +0.92] & \\
 & $\bar r_J$ & +0.91 & [+0.82, +0.96] & [+0.70, +0.95] & \\
F18, compressed & $E_1$ & +0.89 & & [+0.77, +0.96] & \\
 & $E_8$ & +0.96 & & [+0.90, +0.98] & \\
 & $\lambda_1$ & +0.82 & & [+0.49, +0.92] & \\
 & $\bar r_J$ & +0.94 & & [+0.86, +0.99] & \\
X5 & $E_1$ & +0.78 & & [+0.63, +0.84] & +0.79 [+0.69, +0.86] \\
 & $E_8$ & +0.95 & & [+0.91, +0.96] & +0.95 [+0.91, +0.96] \\
 & $\lambda_1$ & +0.67 & & [+0.54, +0.72] & +0.65 [+0.54, +0.72] \\
 & $\bar r_J$ & +0.96 & & [+0.92, +0.97] & +0.97 [+0.95, +0.97] \\
\bottomrule
\end{tabular}}
\end{table}

Without split \{1, 8\}, $\lambda_1$ minus $E_1$ is $-$0.14 [$-$0.26, $-$0.04], $\lambda_1$ minus $E_8$ is
$-$0.30 [$-$0.40, $-$0.22], and $\bar r_J$ minus $E_8$ is +0.02 [+0.00, +0.05].

Table~\ref{tab:app-tests} lists the registered differences. Holm's correction keeps the chance of any false finding among
the 7 secondary tests at 0.05, and 2 of them survive. Its $p$-values come from a
Monte Carlo resampling, so a value close to 0.05 is itself uncertain (``near'' marks the tests within two
standard errors of the threshold).

\begin{table}[h]
\centering
\caption{Registered differences in Spearman correlation (first minus second check; Q1 and Q2 at 97.5\%, the rest at 95\%),
with the registered and the two-level interval, whether the registered verdict still holds, and, for the seven secondary
tests, Holm's adjusted $p$ with whether it survives and whether it is near the threshold. Two-level and Holm columns are exploratory.}
\label{tab:app-tests}
\footnotesize
\setlength{\tabcolsep}{3pt}
\fittable{\begin{tabular}{@{}lccccc@{}}
\toprule
Test & Estimate & Registered & Two-level & Holds & Holm $p$ (survives; near) \\
\midrule
Q1: $\lambda_1-E_1$ at 100 steps, F18 & +0.48 & [+0.20, +0.79] & [+0.09, +1.11] & yes & -- \\
Q2: $\lambda_1-E_8$ at 100 steps, F18 & $-$0.04 & [$-$0.20, +0.14] & [$-$0.24, +0.31] & yes & -- \\
Q1 at 24 steps & +0.36 & [+0.13, +0.60] & [$-$0.03, +0.74] & no & -- \\
Q1 at 368 steps & +0.56 & [+0.31, +0.83] & [+0.19, +1.08] & yes & -- \\
Q2 at 24 steps & $-$0.16 & [$-$0.25, $-$0.08] & [$-$0.36, $-$0.05] & yes & -- \\
Q2 at 368 steps & +0.05 & [$-$0.09, +0.22] & [$-$0.12, +0.34] & yes & -- \\
Q3: $\lambda_1-E_1$, five splits & $-$0.12 & [$-$0.22, $-$0.03] & [$-$0.25, +0.01] & no & 0.056 (no; yes) \\
Q3: $\lambda_1-E_8$, five splits & $-$0.28 & [$-$0.37, $-$0.22] & [$-$0.41, $-$0.22] & yes & $<$0.001 (yes; no) \\
Q4: $\bar r_J-E_1$, F18 & +0.63 & [+0.41, +0.88] & [+0.35, +1.09] & yes & $<$0.001 (yes; no) \\
Q4: $\bar r_J-E_8$, F18 & +0.12 & [+0.02, +0.23] & [+0.03, +0.34] & yes & 0.053 (no; yes) \\
Q5: $\bar r_J-E_1$, compressed & +0.27 & [+0.04, +0.51] & [$-$0.02, +0.74] & no & 0.056 (no; yes) \\
Q5: $\bar r_J-E_8$, compressed & +0.01 & [$-$0.16, +0.20] & [$-$0.22, +0.34] & yes & 0.944 (no; no) \\
Q6: growth-rate shift of pruned copies & +0.03 & [$-$0.02, +0.07] & [$-$0.04, +0.09] & yes & 0.547 (no; no) \\
\bottomrule
\end{tabular}}
\end{table}

\section{Compression and the growth rate}\label{app:A5}
Table~\ref{tab:app-comp} gives, for every compression variant of the 25 five-split parents, the shift in
growth rate, how many copies have $\lambda_1>0$, and the skill retained $K$ (the share of the parent's advantage over
predicting no change that the copy keeps). In all, 7 of the 150 copies with
weights or units removed have $\lambda_1>0$, and 7 of the 75 copies pruned by 50\% or
reduced in width.

\begin{table}[h]
\centering
\caption{Compressed copies of the five-split parents: mean growth rate $\lambda_1$ of the copies, mean shift from the
parent (95\% interval), copies with $\lambda_1>0$, and skill retained $K$ at 8, 100 and 368 steps (negative: worse than predicting
no change).}
\label{tab:app-comp}
\footnotesize
\setlength{\tabcolsep}{3pt}
\fittable{\begin{tabular}{@{}lccccccc@{}}
\toprule
Variant & $\lambda_1$ (1/s) & Shift & Copies $\lambda_1>0$ & $K_8$ & $K_{100}$ & $K_{368}$ \\
\midrule
16-bit weights & $-$0.45 & +0.00 [+0.00, +0.00] & 0 & 100\% & 100\% & 100\% \\
8-bit weights & $-$0.46 & $-$0.01 [$-$0.03, +0.01] & 0 & 99.7\% & 97.0\% & 95\% \\
10\% pruned & $-$0.43 & +0.02 [+0.00, +0.04] & 0 & 99.7\% & 97.1\% & 91\% \\
20\% pruned & $-$0.43 & +0.02 [+0.00, +0.04] & 0 & 97.2\% & 80\% & 57\% \\
30\% pruned & $-$0.42 & +0.03 [$-$0.02, +0.07] & 0 & 90\% & 39\% & $-$41\% \\
50\% pruned & $-$0.19 & +0.27 [+0.23, +0.31] & 0 & 61\% & $-$93\% & $-$445\% \\
Width 192 & $-$0.12 & +0.33 [+0.29, +0.38] & 4 & 30\% & $-$87\% & $-$378\% \\
Width 128 & $-$0.06 & +0.40 [+0.36, +0.44] & 3 & 8\% & $-$182\% & $-$784\% \\
\bottomrule
\end{tabular}}
\end{table}

An earlier registered study of the same recipe, with 500 fine-tuning iterations, is summarised in
Table~\ref{tab:app-perc}. Fine-tuned narrow copies are compared with the equally fine-tuned full-width model, as in the
body; the control row fine-tunes the unreduced model in the same way. Q4 is that study's comparison of fine-tuning
with training from scratch at the same width.

\begin{table}[h]
\centering
\caption{Skill retained $K$ of fine-tuned narrow copies at 100 and 368 steps in the first run, its replication and
the pooled runs (95\% intervals); the unreduced control (pooled); and the registered label and difference for fine-tuning
against training from scratch at 368 steps.}
\label{tab:app-perc}
\footnotesize
\setlength{\tabcolsep}{3pt}
\fittable{\begin{tabular}{@{}llccc@{}}
\toprule
Reading & Width & $K_{100}$ & $K_{368}$ & Q4 (368 steps) \\
\midrule
First run & 192 & 98\% [95\%, 100\%] & 84\% [77\%, 90\%] & not settled [$-$0.06, +0.10] \\
 & 128 & 88\% [82\%, 95\%] & 53\% [40\%, 65\%] & scratch better [$-$0.21, $-$0.05] \\
Replication & 192 & 93\% [89\%, 96\%] & 90\% [84\%, 94\%] & not settled [+0.00, +0.25] \\
 & 128 & 82\% [74\%, 87\%] & 44\% [3.6\%, 65\%] & not settled [$-$0.27, +0.05] \\
Pooled & 192 & 95\% [92\%, 97\%] & 88\% [82\%, 92\%] & not settled [$-$0.01, +0.19] \\
 & 128 & 84\% [77\%, 89\%] & 47\% [19\%, 62\%] & scratch better [$-$0.21, $-$0.01] \\
Control (unreduced) & & 102\% [101\%, 103\%] & 104\% [101\%, 108\%] & -- \\
\bottomrule
\end{tabular}}
\end{table}

\section{Run-time signals}\label{app:A6}
A check that ranks whole simulators need not flag individual bad steps. Earlier analyses of these models, on
the same held-out trajectories of the five splits, found that none of the label-free signals is a usable alarm for a
step that is worse than holding the last state. On the rollout-trained ensembles, a combined kinematic-consistency
score reached an area under the ROC curve (AUROC; 0.5 is chance) of 0.54 [0.40, 0.68], ensemble
disagreement 0.59 and a plain step counter 0.50. Averaged over models of all three kinds, the growth
rate read step by step (over the last eight steps) reached 0.50, below the step counter's 0.64
(difference $-$0.15 [$-$0.22, $-$0.03]; registered: \textsc{worse than counter}). Re-synchronising the simulator with the recorded state when the
kinematic score fired gave 10.1\% [5.6\%, 15.7\%] more 368-step error than the same number of
re-synchronisations on a fixed schedule (registered: \textsc{schedule better}). Physics helps as a repair only where it is exact: rescaling the predicted
gravity direction to unit length changed the 368-step error by $-$3.5\% [$-$5.3\%, $-$2.1\%], while forcing the joint
angles onto the approximate integration rule multiplied it by 2.3.

\section{Study FT: fine-tuning the pruned copies}\label{app:A7}
% A7: the study registered in prereg/PREREG_FT.md before any fine-tuned model existed (R3): design, arms, decision
% rules and results from scripts/ft_stats.py, with any departure recorded in ANOMALIES.md. R9 writes it.
